\chapter{Algoritmo polinomial para conjunto de intervalo e conjunto de envoltória na convexidade P3 orientada}\label{4-algoritmos}

Em~\cite{Medeiros2024Convexidade}, vemos que calcular o número de intervalo e o número de envoltória em árvores orientadas é polinomial, o resultado apresentado segue do \emph{teorema de Courcelle}~\cite{courcelle2012graph}. Nesta seção, apresentamos os algoritmos explícitos dos respectivos parâmetros com estratégias semelhantes para os dois problemas.

Apresentamos o algoritmo~1 o \emph{conjunto de intervalo} e o algoritmo~2 para \emph{conjunto de envoltória} em árvores orientadas. Consideremos $S \subseteq V(T)$ um conjunto de intervalo e que qualquer vértice extremo $v$ obrigatoriamente está em $S$, seja $Ext(T)$ o conjunto de vértices extremos em $T$.

No algoritmo~1, definimos os parâmetros $r^-(v)$ e $r^+(v)$ que são basicamente variáveis inicializadas em 1, para todo $v \in V(T)$ que indicam quantos vizinhos de entrada e saída, respectivamente, são necessários para gerar o vértice $v$.

Já no algoritmo~2, definimos os parâmetros $inf(v)$, $temE(v)$ e $temS(v)$, que recebem valores \emph{booleanos} e que indicam para todo vértice $v \in V(T)$ se $v$ foi gerado, se tem um vértice na vizinhança de entrada que está em $S$ ou está sendo gerado e o mesmo para um vértice na vizinhança de saída, respectivamente.

\begin{algorithm}[ht]
  \SetAlgoLined
  \Entrada{$T, Ext(T)$}
  \Saida{\emph{Conjunto de Intervalo de T}}
  Inicialize $r^-(v) = 1$ e $r^+(v) = 1$ para todo $v \in V(T)$\\
  $S \leftarrow Ext(T)$\\
  Seja $(v_1, v_2, \ldots, v_n)$ uma pós-ordem de $T$\\
  \Para{$i \leftarrow 1, 2, \ldots, n$}{
    Seja $\pi$ pai de $v_i$ em $T$\\
    \Se{$v \in S$}{
      \Se{$\pi \in N^+(v)$}{
        $r^-(\pi) \leftarrow r^-(\pi) - 1$
      }
      \Se{$\pi \in N^-(v)$}{
        $r^+(\pi) \leftarrow r^+(\pi) - 1$
      }      
    }
    \Se{$v \notin S$}{
      \Se{$r^-(v) = 1$ e $r^+(v) = 1$}{
        $S \leftarrow S \cup \{v\}$
      }
      \Se{$r^-(v) \leq 0$ e $r^+(v) = 1$}{
        \Se{$\pi \in N^+(v)$ e $\pi \notin S$}{
          $S \leftarrow S \cup \{\pi\}$
        }
        \Se{$\pi \in N^-(v)$}{
          $S \leftarrow S \cup \{v\}$\\
          $r^+(\pi) \leftarrow r^+(\pi) - 1$
        }
      }
      \Se{$r^-(v) = 1$ e $r^+(v) \leq 0$}{
        \Se{$\pi \in N^-(v)$ e $\pi \notin S$}{
          $S \leftarrow S \cup \{\pi\}$
        }
        \Se{$\pi \in N^+(v)$}{
          $S \leftarrow S \cup \{v\}$\\
          $r^-(\pi) \leftarrow r^-(\pi) - 1$
        }
      }
    }
  }
  \Retorna{$S$}
  \caption{Conjunto de Intervalo em Árvores orientadas}\label{alg1}
\end{algorithm}

\begin{theorem}
  Algoritmo\ref{alg1} retorna o menor conjunto de intervalo para a árvore T
  \begin{proof}
  Seja $I^* = (i^*_1, \ldots, i^*_j, \ldots, i^*_n)$ uma sequência de escolhas ótimas para obtermos o \emph{conjunto de intervalo} de T, onde $i^*_i = 1$ se o vértice $v_i$ está no \emph{conjunto de intervalo} e $0$ caso contrário. Além disso, seja $S = (s_1, \ldots, s_j, \ldots, s_n)$ a sequência de escolhas feita pelo \emph{Algoritmo 1}. Se $I^*=S$, está resolvido. Suponha que $I^* \neq S$, neste caso, seja $i$ a primeira posição em que as sequências diferem.

  Suponha primeiro que $i^*_i = 0$ e $s_i = 1$, ou seja, o vértice $v_i$ foi escolhido pelo \emph{algoritmo} mas não pela sequência $I^*$. Com isso temos os seguintes casos:

  \begin{enumerate}
    \item Neste caso, $v_i$ entrou na condicional \textbf{15}, como $I^*$ e $S$ concordam em todas as posições anteriores a $i$, temos que os filhos de $v_i$ não estão em $I^*$, caso contrário, o algoritmo não entraria na condicional 15. O único vizinho restante de $v_i$ é seu pai $\pi$, mas mesmo que $\pi \in I^*$, se $v_i \notin I^*$, então $v_i$ não seria gerado por $I^*$, contradizendo o fato de $I^*$ ser \emph{Interval Set}.
    \item Neste caso, $v_i$ entrou na condicional \textbf{22}, como $I^*$ e $S$ concordam em todos os filhos de $v_i$, então $v_i$ possui pelo menos um filho em $N^-(v_i)$ que está em $S$, mas nenhum filho em $N^+(v_i)$ que está em $S$. Além disso, o pai $\pi \in N^-(v_i)$. O único vizinho restante de $v_i$ é seu pai $\pi$, mas mesmo que $I^*$ escolha $\pi$, como $\pi \in N^-(v_i)$ e $I^*\cap N^+(v_i) = \emptyset$, então $I^*$ não gera $v_i$, um absurdo.
    \item Condicional \textbf{31} é análoga ao caso anterior.
  \end{enumerate}

  Agora suponha que $i^*_i=1$ e $s_i = 0$, ou seja, o vértice $v_i$ foi escolhido pela sequência $I^*$ mas não pelo \emph{algoritmo}. Com isso temos os seguintes casos:

  \begin{enumerate}
    \item O algoritmo não escolheu $v_i$ pois entrou na condicional \textbf{19}. Como $v_i \in I^*$, então seu pai $\pi \notin I^*$, pois caso contrário, obteríamos um conjunto de intervalo menor, neste caso, $I^*\backslash \{v_i\}$. Mas perceba que $I'^* = I^*\backslash \{v_i\} \cup \{\pi\}$ continua sendo conjunto de intervalo tal que $|I^*| = |I'^*|$, portanto, mínimo e concorda em mais uma posição com a sequência $S$.
    \item Condicional \textbf{28} é análoga ao caso anterior.
  \end{enumerate}
  \end{proof}
\end{theorem}

\begin{algorithm}[ht]
  \SetAlgoLined
  \Entrada{$T, Ext(T)$}
  \Saida{\emph{Conjunto de Envoltória de T}}
  $\forall v \in V(T)$: $inf(v) \leftarrow false$, $temE(v) \leftarrow false$ e $temS(v) \leftarrow false$\\
  $S \leftarrow Ext(T)$\\
  $\forall v \in S$: $inf(v) \leftarrow true$\\
  Seja $(v_1, v_2, \ldots, v_n)$ uma pós-ordem de $T$\\
  \Para{$i \leftarrow 1, 2, \ldots, n$}{
    Seja $\pi$ pai de $v_i$ em $T$\\
    \Se{$\exists u \in (N^-(v_i) \cap filhos(v_i))$ \emph{tal que} $inf(u)=true$}{
      $temE(v_i) \leftarrow true$      
    }
    \Se{$\exists u \in (N^+(v_i) \cap filhos(v_i))$ \emph{tal que} $inf(u)=true$}{
      $temS(v_i) \leftarrow true$      
    }
    \Se{$temE(v_i) = true$ \emph{e} $temS(v_i) = true$}{
      $inf(v_i) \leftarrow true$      
    }
    \Se{$v_i \notin S$ \emph{e} $inf(v_i) = false$}{
      \Se{$\pi = NULO$ \emph{ou} $(temE(v_i) = false$ \emph{e} $temS(v_i) = false)$}{
        $S \leftarrow S \cup \{v_i\}$\\
        $inf(v_i) \leftarrow true$
      }
      \SenaoSe{$temE(v_i) = true$ \emph{e} $temS(v_i) = false$ \emph{e} $\pi \in N^-(v_i)$}{
        $S \leftarrow S \cup \{v_i\}$\\
        $inf(v_i) \leftarrow true$
      }
      \SenaoSe{$temE(v_i) = false$ \emph{e} $temS(v_i) = true$ \emph{e} $\pi \in N^+(v_i)$}{
        $S \leftarrow S \cup \{v_i\}$\\
        $inf(v_i) \leftarrow true$
      }
    }
  }
  \Retorna{$S$}
  \caption{Conjunto de envoltória em Árvores orientadas}\label{alg2}
\end{algorithm}

\begin{theorem}
  O Algoritmo 2 retorna o menor conjunto de envoltória para a árvore T
\end{theorem}

\begin{proof}
  Seja $I^* = (i^*_1, \ldots, i^*_j, i^*_n)$ uma sequência de escolhas ótimas para obtermos um \emph{conjunto de envoltória} de $T$, onde $i^*_i = 1$, se o vértice $v_i$ está no \emph{conjunto de envoltória} e 0 caso contrário. Além disso, seja $S = (s_1, \ldots, s_j, \ldots, s_n)$ uma sequência de escolhas feita pelo \emph{Algoritmo 2}. Se $I^* = S$, está resolvido. Suponha que $I^* \neq S$, neste caso, seja $i$ a primeira posição em que as sequências diferem.

  Suponha primeiro que $i^*_i = 0$ e $s_i = 1$, ou seja, $v_i$ foi escolhido pelo algoritmo, mas não foi escolhido pela sequência $I^*$. Com isso temos os seguintes casos:

  \begin{enumerate}
    \item Neste caso, $v_i$ pode ter entrado na condicional \textbf{17}, oberve que, ou $v_i$ é a raiz (não tem pai $\pi$), ou não é gerado pelos seus filhos, pois, temos que $inf(v_i) = false$. Como $I^*$ e $S$ concordam em todas as posições anteriores a $i$, temos que neste caso, os filhos de $v_i$ não estão em $I^*$, caso contrário, não entrária nessa condicional, já que dada a condição \textbf{13}, $v_i$ seria gerado. Como $\pi = NULO$ e $v_i$ não é gerado por seus filhos, a única opção para $v_i$ é ser adicionado ao conjunto de envoltória. Como $I^*$ não selecionou $v_i$, não temos um conjunto de envoltória, já que $v_i$ não pode ser gerado. Contradição.
    \item A outra possibilidade é $v_i$ ter entrado na condicional \textbf{21}, neste caso, $v_i$ só tem o filho de entrada que ajuda a gerá-lo e além disso, não importa o estado de $\pi$, visto que $\pi \in N^-(v_i)$, não contribuindo para gerar $v_i$. A única opção que resta para $v_i$ é o mesmo estar no conjunto de envoltória. Portanto, se $v_i$ não foi selecionado por $I^*$, não temos um conjunto de envoltória. contradição.
    \item A possibilidade da condicional \textbf{25} é análoga ao caso anterior.
  \end{enumerate}

  Agora suponha que $i^*_i = 1$ e $s_i = 0$, ou seja, $v_i$ foi escolhido por $I^*$, mas não foi escolhido pela pelo algoritmo. Com isso temos os seguintes casos:

  \begin{enumerate}
    \item O algoritmo não selecionou $v_i$ pois, na linha \textbf{14}, temos que $inf(v_i) = true$. Neste caso, temos que $v_i \in I^*$, mas perceba que como $v_i$ é gerado por seus filhos, $I^*\backslash \{v_i\}$ é um conjunto de envoltória menor que $I^*$. Contradição.
    \item O algoritmo não entrou nas condições \textbf{17}, \textbf{21} e \textbf{25}. Neste caso, se $\pi \neq NULO$, como $inf(v_i) = false$, temos pelo menos um filho $u$ de $v_i$ de entrada ou de saída tal que $inf(u) = true$. Suponha que $temE(v_i) = true$ e $temS(v_i) = false$. Como $v_i$ não entrou na condicional \textbf{21}, temos que $\pi \in N^+(v_i)$. Temos os seguintes casos:
    \begin{enumerate}
      \item Se $\pi \in I^*$, então teríamos que $I^*\backslash \{v_i\}$ seria um conjunto de intervalo menor, contradizendo a otimalidade de $I^*$.
      \item Agora, se $\pi \notin I^*$, perceba que $I'^* = I^*\backslash \{v_i\} \cup \{\pi\}$, continua sendo um conjunto de envoltória e que $|I'^*| = |I^*|$, portanto, mínimo e que como $v_i \notin I'^*$, concorda em mais um posição com o algoritmo, contradizendo a escolha de $I^*$.
    \end{enumerate}

    \item Para a condição \textbf{25}, é análogo a análise anterior.
  \end{enumerate}
\end{proof}
